\week{10}
\lecture{25}{continuation of Integrals, i guess...}
% \section*{Week 10. Lecture 25}

We have started to study the complex integration, formulated and proved Prop 17.3 about contour integrals. We conclude the following useful corollary.
\begin{corollary}\label{corollary:17.4}
    Let $f$ be continuous on a domain D, with antiderivative F on D.
    
    1) If $C_1$ and $C_2$ are any two contours both starting at $z_0$ and both ending at $z_1$, then $\int_{C_1} f(z) dz = \int_{C_2} f(z) dz$.
    
    2) If C is a closed contour in D, then $\int_C f(z) dz = 0$.
\end{corollary}
\begin{proof}
    1) Immediate from Prop 17.3 4)
    
    2) If C is a closed contour, then we can write $C = C_1 - C_2$, where $C_1$ and $C_2$ have the same start point and the same end point. Then
    \[  \int_C f(z) dz = \int_{C_1 - C_2} f(z) dz = \int_{C_1} f(z) dz + \int_{-C_2} f(z) dz = \int_{C_1} f(z) dz - \int_{C_2} f(z) dz = 0 \]
    Prop 17.3 1)
    
\end{proof}
Note: The condition that $f$ has antiderivative is necessary!

\begin{example}\label{example:19.11} Compute $\int_C f(z) dz$ where $f(z) = \bar{z}$ and C connects 2 points $z_0 = 0, z_1 = 2 + 4i$

1) by straight line segment
2) by parabola (part of it) $y = x^2$

Sol: 1) Straight line connecting 0 and $2 + 4i$ is parametrized by $\gamma(t) = (2+4i)t, 0 \leq t \leq 1$. Thus, $\int_C f(z) dz = \int_0^1 \overline{f(\gamma(t))} \gamma'(t) dt = \int_0^1 (2+4i) dt = \int_0^1 20t dt = 10$.

2) Parametrization: $\gamma(t) = t + it^2, 0 \leq t \leq 2$, and $\int_C f(z) dz = \int_0^2 \overline{f(\gamma(t))} \gamma'(t) dt = \int_0^2 (t-it^2)(1+2it) dt = \int_0^2 (t-it^2)(1+2it) dt = \int_0^2 (t + it^2 + 2t^3) dt = 10 + i\frac{8}{3} (\neq 10)$.

$\bar{z}$ does not have an antiderivative!
\end{example}

\begin{example}\label{example:19.12} Let $C_1$ be the contour from -1 to 1 along the upper unit semicircle and $C_2$ the contour from -1 to 1 along the 
lower unit semicircle. Compute $\int_{C_1} \frac{1}{(z-4)^2} dz$.

Sol: Note: $\int_{C_1 - C_2} \frac{1}{(z-4)^2} dz = 0$. This is an analytic function on the unit circle, the only singularity is a 
pole of order 2 at $z = 4$ - outside the unit circle, with antiderivative $-\frac{1}{z-4}$.

$\implies$ By Cor 17.4 2(): 
\[ \int_{C_1-C_2}^{} \frac{1}{(z-4)^2} dz = 0.\]

 $C_1 - C_2$ is a closed contour; the unit circle.
\[ \implies \int_{C_1} \frac{dz}{(z-4)^2} = \int_{C_2} \frac{dz}{(z-4)^2} = -\frac{1}{z-4}|_{-1}^1 = \frac{1}{-3} - \frac{1}{-5} = \frac{2}{15} .\]
\end{example}

\begin{example} \label{definition:19.13} Compute $\int_C \frac{dz}{(z-z_0)^n}$, where $C = \{z \in \mathbb{C}: |z-z_0| = R\}$ and $n$ is some integer.

Sol: C is a circle centered at $z_0$ of radius R. Parametrization of C: $|z-z_0| = R \implies (x-x_0)^2 + (y-y_0)^2 = R^2$. 
Take $x-x_0 = R\cos t, y - y_0 = R\sin t, 0 \leq t \leq 2\pi$, then by Euler's formula $z - z_0 = R(\cos t + i\sin t) = Re^{it}, 0 \leq t \leq 2\pi \implies dz = iRe^{it} dt$.
\end{example}

By Def 17.7 (of the contour integral) $\int_C \frac{dz}{(z-z_0)^n} = \int_0^{2\pi} \frac{iRe^{it}}{R^n e^{int}} dt = \frac{i}{R^{n-1}} \int_0^{2\pi} e^{i(1-n)t} dt$.

Consider two cases:

$n = 1: \int_C \frac{dz}{z-z_0} = i \int_0^{2\pi} dt = 2\pi i$.

$n \neq 1: \int_C \frac{dz}{(z-z_0)^n} = \frac{i}{R^{n-1}} \int_0^{2\pi} e^{i(1-n)t} dt = \frac{i}{(1-n)R^{n-1}} e^{i(1-n)t} |_0^{2\pi} = \frac{1}{(1-n)R^{n-1}} [e^{i(1-n)2\pi} - 1] = 0$ \{$e^{2\pi ki} = 1 \quad \forall k \in \mathbb{Z}$\}.

\begin{example}\label{example:19.4} Compute $\int_C \frac{dz}{\sqrt{z}}$ where $C = \{f(x,y): x^2 + y^2 = 16, y \geq 0\}$, and the branch of $\sqrt{z}$ is such that $\sqrt{4} = -2$.

Sol: Note: the contour is the upper half of the circle centered at 0 of radius 4.

Rewrite: $C = \{z \in \mathbb{C}: |z| = 4, \text{Im } z > 0\} \implies$ if $z \in \mathbb{C}$, then $z = 4e^{it}$ for $0 \leq t \leq \pi$.

$\sqrt{z}$ attains 2 values: $w_1 = \sqrt{4e^{it}} = 2e^{i\frac{t}{2}}$ and $w_2 = \sqrt{4e^{it}} = 2e^{i(\frac{t}{2} + \pi)}$.

If $t = 0: w_1 = \sqrt{4} = 2e^0 = 2$ and $w_2 = \sqrt{4} = 2e^{i\pi} = -2 \implies$ the given branch is $w_2$. Now we compute: $\sqrt{z} = 2e^{i(\frac{t}{2} + \pi)}, dz = 4ie^{it} dt \implies \int_C \frac{dz}{\sqrt{z}} = \int_0^\pi \frac{4ie^{it}}{2e^{i(\frac{t}{2} + \pi)}} dt = 2i \int_0^\pi e^{i(\frac{t}{2} - \pi)} dt = 4e^{i(\frac{t}{2} - \pi)}|_0^\pi = 4(e^{-i\frac{\pi}{2}} - e^{-i\pi}) = 4(-i + 1)$ \{$\because e^{-i\frac{\pi}{2}} = \cos \frac{\pi}{2} - i\sin \frac{\pi}{2} = -i$\}.
\end{example}

The estimate in Prop 17.2 for the size of an integral of a complex-valued function of a real variable also has an important consequence for contour integrals.
\begin{proposition}[ML Inequality]\label{proposition:17.5} Let C be a contour. If $|f(z)| \leq M$ for all $z \in C$ and Length (C) = L, then $|\int_C f(z) dz| \leq ML$.
\end{proposition}
\begin{proof}
    $|\int_C f(z) dz| = |\int_a^b f(\gamma(t)) \gamma'(t) dt| \leq \int_a^b |f(\gamma(t))| |\gamma'(t)| dt$
    $\leq M \int_a^b |\gamma'(t)| dt \overset{\text{Def 17.5}}{=} ML$.
\end{proof}

Now that we have defined contour integrals and discussed some of their properties, we concentrate on one of the key theorems of complex analysis - Cauchy's Theorem.

\textbf{Cauchy's Theorem}

\begin{definition} \label{definition:18.1} $U \subseteq \mathbb{C}$ is convex if for all $z_1, z_2 \in U$, the line segment $[z_1, z_2] = \{tz_2 + (1-t)z_1: 0 \leq t \leq 1\}$ (is contained in U).
    
$U \subseteq \mathbb{C}$ is star-shaped about $w$ if for all $z \in U$, the line segment $[z, w] \subseteq U$.
\end{definition}

\begin{remark}
    Any convex set is star-shaped, but not every star-shaped set is convex: look at  - it is star-shaped about $w$ (in the center), but it is not convex - the line segment connecting any two neighboring corners is not contained in this set.
\end{remark}

\begin{remark}
    Set is convex if and only if it is star-shaped with respect to any point in that set. Note: star-shaped set is always path-connected, thus it is always connected.
\end{remark}

Our goal now is to prove Cauchy's Theorem for star-shaped domains.
\begin{theorem}[Cauchy's Thm for star-shaped domains]
    Let $f$ be a function which is analytic on an open star-shaped region $U \subseteq \mathbb{C}$. Then, for every closed contour C in U $\int_C f(z) dz = 0$.
\end{theorem}
\begin{proof}
    
    It will suffice to show that $f$ has an antiderivative F on U, since then by Cor 17.4 2) we will conclude for a closed contour C that $\int_C f(z) dz = 0$.
    
    Define: $F(z) = \int_{[w, z]} f(\xi) d\xi$, where w is such that U is star-shaped about $w$, and $[w, z]$ is the straight line segment from $w$ to $z$.
    
    Claim: F is differentiable and $F'(z) = f(z)$.
    
    Pf of Claim: Consider $\frac{F(z) - F(z_0)}{z - z_0} = \frac{1}{z-z_0} [\int_{[w, z]} f(\xi) d\xi - \int_{[w, z_0]} f(\xi) d\xi]$.
    
    Suppose we could show that (*) $\int_{[w, z]} f(\xi) d\xi - \int_{[w, z_0]} f(\xi) d\xi = \int_{[z_0, z]} f(\xi) d\xi$.
    
    Then, by (*): $\frac{F(z) - F(z_0)}{z - z_0} = \frac{1}{z-z_0} \int_{[z_0, z]} f(\xi) d\xi$.
    
    $f$ is analytic on $U \implies$ by Prop 6.5 $f$ is continuous on $U \implies$ for every $\epsilon > 0$ there exists $\delta > 0$ s.t. if $|z - z_0| < \delta$, then $|f(z) - f(z_0)| < \epsilon$.
    $\implies$ If $|z - z_0| < \delta$, then $|\int_{[z_0, z]} f(\xi) d\xi - \int_{[z_0, z]} f(z_0) d\xi| < \epsilon |z - z_0|$, 
    since $|\int_{[z_0, z]} f(\xi) d\xi - \int_{[z_0, z]} f(z_0) d\xi| \leq \epsilon |z - z_0|$.
    
    Thus: $|\frac{F(z) - F(z_0)}{z - z_0} - f(z_0)| < \epsilon \implies$ F is differentiable at $z_0$, 
    with derivative $F'(z_0) = f(z_0) \implies$ We proved Cauchy's Theorem up to (*). It remains to prove (*) which is Cauchy's Theorem for a Triangle.
\end{proof}

% \section*{Lectures 26 + 27.}
\lecture{26+27}{Cauchy theorem for a Triangle, and other stuff...}

We have proved Cauchy's Theorem for star-shaped domains up to: (*) $\int_{[w, z]} f(\xi) d\xi - \int_{[w, z_0]} f(\xi) d\xi = \int_{[z_0, z]} f(\xi) d\xi$, which is Cauchy's Theorem for a Triangle. Now we prove it.
\begin{theorem}[Cauchy's Theorem for a triangle]
    Let $f$ be analytic in a domain U containing a triangle T. Then: $\int_{\partial T} f(z) dz = 0$.
\end{theorem}
\begin{proof}
    
    Let $\eta(T) = \int_{\partial T} f(z) dz$, denote the length of $\partial T$ by L.
    
    Divide T into 4 triangles $T^{(j)}$ by bisecting the sides of T.
    
    By orienting each of the boundaries $\partial T^{(j)}$ of the subtriangles anticlockwise, 
we have: $\eta(T) = \sum_{j=1}^4 \eta(T^{(j)})$ \{the internal edges of each of the subtriangles $T^{(j)}$ will cancel out\}. 
Each $T^{(j)}$ has a boundary with Length $(\partial T^{(j)}) = \frac{L}{2}$. At least one of the $T^{(j)}$ must satisfy $|\eta(T^{(j)})| \geq \frac{1}{4} |\eta(T)|$. Call this triangle $T_1$.
    
    Divide $T_1$ into 4 triangles in the same way as T and conclude that one of these triangles, $T_2$, satisfies $|\eta(T_2)| \geq \frac{1}{4} |\eta(T_1)| \geq \frac{1}{4^2} |\eta(T)|$, 
Length $(\partial T_2) = \frac{L}{4}$. Repeating this procedure, we construct a sequence of 
triangles $T_1 \supset T_2 \supset T_3 \supset \dots \supset T_n \supset \dots$ with $|\eta(T_n)| \geq \frac{1}{4^n} |\eta(T)|$ and Length $(\partial T_n) = \frac{L}{2^n}$. Let $z_0$ be the point of intersection $\bigcap_{n=1}^\infty T_n$ \{exists and unique by Cantor Lemma or (**)\}. 
Since $f$ is differentiable at $z_0$, given any $\epsilon > 0$ there exists $\delta > 0$ such that for $|z - z_0| < \delta$ we have $|f(z) - f(z_0) - f'(z_0)(z-z_0)| < \epsilon |z - z_0|$.

Take $n$ such that $\frac{L}{2^n} < \delta$. Then $\epsilon |z - z_0| < \frac{L\epsilon}{2^n}$ for every $z \in \partial T_n$, so 
$|\int_{\partial T_n} [f(z) - f(z_0) - f'(z_0)(z-z_0)] dz| \leq \frac{L\epsilon}{2^n} \cdot \frac{L}{2^n} = \frac{L^2}{4^n} \epsilon$.

However: $\int_{\partial T_n} [f(z_0) + f'(z_0)(z-z_0)] dz = 0$, since $f(z_0) + f'(z_0)(z-z_0)$ has antiderivative $z f(z_0) + f'(z_0)(\frac{z^2}{2} - z z_0)$ and $\partial T_n$ is a closed 
curve (Cor 17.4). Hence: $|\eta(T_n)| = |\int_{\partial T_n} f(z) dz| \leq (\frac{L^2}{4^n}) \epsilon$.

On the other hand we KNOW that $|\eta(T_n)| \geq \frac{1}{4^n} |\eta(T)|$
$\implies |\eta(T)| \leq \epsilon L^2$. Since $\epsilon > 0$ is arbitrary $\implies \eta(T) = 0$.

(**) Existence of $z_0$: On each step of the process choose $z_n \in T_n$ (a point) and obtain an infinite sequence $\{z_n\}_{n=1}^\infty$. Then:

1) $\{z_n\}$ is a Cauchy sequence, namely, for any $\epsilon > 0$ there exists $N \in \mathbb{N}$ such that for any pair $n, m > N$, $|z_n - z_m| < \epsilon$.

First, it is bounded. Set $\epsilon = 1$ and choose N such that for any $n, m \in \mathbb{N}$, $|z_n - z_m| < 1$ (using that $\{z_n\}$ is a Cauchy sequence). 
Namely: (*) $||z_n| - |z_m|| < |z_n - z_m| < 1$. Fix $m_0 \in \mathbb{N}$. 
From (*) for all $n > N$, $|z_{m_0}| - 1 < |z_n| < |z_{m_0}| + 1 \implies$ for all $n \geq 1$, $\min\{|z_1|, \dots, |z_{m_0}-1|, |z_{m_0}|+1\} \leq |z_n| \leq \max\{|z_1|, \dots, |z_{m_0}-1|, |z_{m_0}|+1\}$ $\implies$ it is bounded. Now apply to $\{z_n\}$ Bolzano-Weierstrass Thm and obtain a converging subsequence $\{z_{n_k}\}, \lim_{k \to \infty} z_{n_k} = z_0$ (apply separately to Re $z_n$ and Im $z_n$).

2) $\lim z_n = z_0$: Fix $\epsilon > 0$, choose $N_1$ s.t. $\forall n, m > N_1, |z_n - z_m| < \frac{\epsilon}{2}$ (Cauchy property). 
Choose $N_2$ so that $\forall n_k > N_2$ and such that $z_{n_k}$ belongs to the subsequence. We have: $|z_0 - z_{n_k}| < \frac{\epsilon}{2}$ (use that $\lim_{k \to \infty} z_{n_k} = z_0$). Set $N = \max(N_1, N_2)$ and fix $n_k > N$ such that $z_{n_k}$ belongs to the subsequence. Then, for any $n > N$, $|z_0 - z_n| = |z_0 - z_{n_k} + z_{n_k} - z_n| \leq |z_{n_k} - z_0| + |z_n - z_{n_k}| < \frac{\epsilon}{2} + \frac{\epsilon}{2} = \epsilon$ $\implies \lim_{n \to \infty} z_n = z_0$.

3) Let us show that $\{z_n\}$ is a Cauchy sequence: $z_1 \in T_1, z_2 \in T_2, \dots$ 
Since the diameter of each triangle shrinks by a factor 2 (at each generation) and since diam $(T_n) \leq$ Length $(\partial T_n)$ we get for sufficiently large $n: |z_n - z_m| <$ Length $(\partial T_n) < \frac{L}{2^n} < \epsilon \implies \{z_n\}$ is a Cauchy sequence that converges to $z_0$.

\end{proof}
From the statement of Cauchy's Theorem for a star-shaped region, one can deduce the same result for more general regions comprised of star-shaped pieces.

From Thm 18.1 follows: If a curve C can be expressed as $C = C_1 + C_2 + \dots + C_N$ where the region enclosed by each $C_j$ is star-shaped, then
$\int_C f(z) dz = \int_{C_1} f(z) dz + \int_{C_2} f(z) dz + \dots + \int_{C_N} f(z) dz = 0$.

As a consequence one can prove:
\begin{theorem}[Cauchy's Theorem]\label{theorem:18.3} If C is any simple, closed contour in $\mathbb{C}$ and $f$ is analytic everywhere on and inside C, then $\int_C f(z) dz = 0$.
\end{theorem}

One of very useful corollaries of this theorem is the Deformation Principle.
\begin{corollary}[Deformation Principle]\label{corollary:18.4}
    Let $C_1$ and $C_2$ be two simple, closed, positively oriented contours such that $C_1$ lies interior to $C_2$. If $f$ is analytic in a domain D that contains both $C_1$ and $C_2$ and the region between them, then $\int_{C_1} f(z) dz = \int_{C_2} f(z) dz$.
\end{corollary}

The picture that you should have in mind here is where D is a domain that contains both curves and the region between them: [PICTURE!]
\begin{proof}
    The main idea of the proof is to divide the region between $C_1$ and $C_2$ by "cuts" into star-shaped regions.
    
    Construct 2 disjoint contours (or cuts) $L_1$ and $L_2$ that join $C_1$ and $C_2$. The contour $C_1$ is cut into 2 contours $C_1^*$ and $C_1^{**}$ and the contour $C_2$ is cut into $C_2^*$ and $C_2^{**}$.
    
    Now form 2 new contours: $K_1 = -C_1^* + L_1 + C_2^* - L_2$ \{-$L_2, -C_1^*$ since we have changed the orientation\}, $K_2 = -C_1^{**} + L_2 + C_2^{**} - L_1$.
    
    $K_{1,2}$ - new contours, $D_{1,2}$ - new domains.
    
    Now we can apply Thm 18.3 (Cauchy's Thm) to the contours $K_1$ and $K_2$ and obtain: $\int_{K_1} f(z) dz = 0 \quad \int_{K_2} f(z) dz = 0$.
    
    Adding contours (to reconstruct the original one) gives us: $K_1 + K_2 = -C_1^* + L_1 + C_2^* - L_2 - C_1^{**} + L_2 + C_2^{**} - L_1 = C_2^* + C_2^{**} - (C_1^* + C_1^{**}) = C_2 - C_1$.
    
    $\implies 0 = \int_{K_1} f(z) dz + \int_{K_2} f(z) dz = \int_{K_1 + K_2} f(z) dz = \int_{C_2 - C_1} f(z) dz = \int_{C_2} f(z) dz - \int_{C_1} f(z) dz$.
    
    Namely: $\int_{C_2} f(z) dz = \int_{C_1} f(z) dz$.
\end{proof}

Now we can give an alternative proof to Ex 19.13. It is an important tool for computations.
\begin{corollary}\label{corollary:18.5} Let $z_0 \in \mathbb{C}$ be fixed. If C is a simple, closed, positively oriented contour such that $z_0$ lies interior to C, then $\int_C \frac{dz}{(z-z_0)^n} = \begin{cases} 2\pi i & n = 1 \\ 0 & n \neq 1, n \in \mathbb{Z} \end{cases}$.
\end{corollary}
\begin{proof}
    Since $z_0$ lies interior to C, we can choose R so that the circle $C_R$ with center $z_0$ and radius R lies interior to C.
        Hence, $f(z) = \frac{1}{(z-z_0)^n}$ is analytic in a domain D that contains both C and $C_R$ and the region between them \{we excluded the only singularity of $f$ $z = z_0$, everywhere else $f$ is analytic\}.
    
    As in Ex 19.13: let $C_R$ be parametrized by: $z(\theta) = z_0 + Re^{i\theta}, 0 \leq \theta \leq 2\pi$, $dz(\theta) = iRe^{i\theta} d\theta$.
    
    The deformation principle implies that $\int_C \frac{dz}{(z-z_0)^n} = \int_{C_R} \frac{f(z) dz}{(z-z_0)^n} = \int_0^{2\pi} \frac{iRe^{i\theta}}{R^n e^{in\theta}} d\theta = \frac{i}{R^{n-1}} \int_0^{2\pi} e^{i(1-n)\theta} d\theta = \begin{cases} i \int_0^{2\pi} d\theta = 2\pi i & n = 1 \\ \frac{1}{(1-n)R^{n-1}} e^{i(1-n)\theta} |_0^{2\pi} = \frac{1}{(1-n)R^{n-1}} - \frac{1}{(1-n)R^{n-1}} = 0 & n \neq 1 \end{cases}$.
\end{proof}

More generally, a similar proof to the one used for the Deformation Principle, gives
\begin{corollary}\label{corollary:18.6}
    If C is a simple, closed, positively oriented contour, and $C_1, C_2, \dots, C_n$ are simple, closed, positively oriented contours inside C with disjoint interiors, and $f$ is analytic in the region between C and $C_1, C_2, \dots, C_n$, then $\int_C f(z) dz = \int_{C_1} f(z) dz + \int_{C_2} f(z) dz + \dots + \int_{C_n} f(z) dz$.
\end{corollary}

The picture to have in mind is:
[PUT PICTRE HERE]

\begin{example}\label{example:20.1} Compute $\int_C \frac{z}{(z-4)^2} dz$, where C is the contour defined by traversing once the square with vertices $\pm 3i, 6 \pm 3i$ anticlockwise.

Sol: To compute this integral we apply the Deformation principle and observe that $\int_C \frac{z dz}{(z-4)^2} = \int_{C'} \frac{z dz}{(z-4)^2}$, where C' is the circular contour centered at 4 of radius 1. In particular, C' may be parametrized by the path $\gamma(t) = 4 + e^{2\pi i t}, 0 \leq t \leq 1$
$\implies \int_{C'} \frac{z dz}{(z-4)^2} = \int_0^1 \frac{4 + e^{2\pi i t}}{e^{4\pi i t}} 2\pi i e^{2\pi i t} dt = 2\pi i \int_0^1 (4e^{-2\pi i t} + 1) dt = 2\pi i [-\frac{2}{\pi i}e^{-2\pi i t} + t]|_0^1 = [-4e^{-2\pi i t} + 2\pi i t]|_0^1 = -4 + 2\pi i - (-4 + 0) = 2\pi i$.

\end{example}

